% This file should be rename with the speaker's last name. (e.g : smith)

% Write the author names in the order you wish them to appear
% and underline the one who will give the talk.

% Only the speaker's email address is required.

\begin{abstract_online}{Computer Algebra, Student Assessment and Learning Data Analysis}{%
        David Smith$^1$, \underline{Stephen M. Watt}$^{2}$}{[smwatt@uwaterloo.ca]
        }{%
        $^1$ Digital Education Company Ltd, Cambridge, UK\newline{}
        $^2$ David R. Cheriton School of Computer Science, University of Waterloo, Canada}}
    
Online teaching and assessment tools are typically lacking for mathematical subjects in two different ways:  The first is in producing sufficiently numerous questions of equivalent difficulty for a given topic.  The second is to be able to evaluate student-generated answers that may have many mathematically equivalent forms.  A computer algebra system is needed to do both these things well.   The M\"{o}bius platform uses an embedded Maple system for this purpose.   M\"{o}bius has been deployed to more than 400 educational institutions and has evaluated more than 35,000,000 student problems.  This has been used to collect anonymized data for analysis. This data can be used to measure and identify factors in learner engagement and course pathways. We summarize one study of how this has been used to measure the relationship between uniformity of engagement and student outcomes [2].  The talk concludes with some forward-looking ideas on learner modelling.

\textbf{Conflict of Interest Statement}\newline{}The authors are related to Digital Education Company, Ltd, the M\"{o}bius service provider .

\textbf{Keywords} \newline{} online education, mathematics, teaching and assessment

\textbf{References} \newline{}
[1] \textsc{M\"{o}bius}, \texttt{https://www.digitaled.com/mobius} (retrieved 2022-08-04)\newline{}
[2] \textsc{David Smith, Aron Pasieka, Ralf Becker, Christina Perdikoulias}, 
Student Success in Asynchronous STEM Education:
measuring and identifying contributors to learner.
In \emph{2022 IEEE Global Engineering Education Conference (EDUCON)}, 473-479.\newline{}
\end{abstract_online}
    
